%%%%%%%%%% source scan idx 53 | folio 46 %%%%%%%%%%

1.4. 设 $\underline E$ 为概形 $S$ 上秩为 $r+1$ 的局部自由层, $\mathbb P(\underline E)$ 为相应的射影丛, $d\leq r$ 为整数, $\underline a=(a_1,\ldots,a_d)$ 为整数序列. $\mathbb P(\underline E)$ 中多重次数为 $\underline a$ 的相对\underline{完全交}, 是指在每个纤维中余维均为 $d$ 的闭子概形, 且在 $S$ 的平展拓扑下局部地是次数分别为 $a_1,\ldots,a_d$ 的 $d$ 个超曲面的交.

由 1.1, 对于 $X/S=\mathbb P(\underline E)/S$, $\underline F=\underline O_X$,

\begin{equation*}
\underline W=\bigoplus_i\underline O(-a_i),
\end{equation*}

以及定义多重次数为 $\underline a$ 的相对完全交 $Y\subset\mathbb P(\underline E)$ 的 $\underline W^\vee$ 的截面 $s$, 1.3 的假设 (b) 和 (c) 成立.

\underline{定理} 1.5. \underline{在 1.4 的假设下, 设 $n=r-d$ 为 $Y$ 的相对维数, $\eta\in H^0(S,R^1f_*\Omega^1_{Y/S})$ 为下图中 $\underline O(1)$ 的第一 Chern 类}

\begin{center}
\begin{tikzcd}[column sep=large,row sep=large]
Y\arrow[rr,hook]\arrow[dr,dash,"f"']&&
\mathbb P(\underline E)\arrow[dl,"g"]\\
&S&{}
\end{tikzcd}
\qquad .
\end{center}

\underline{若 $Y$ 在 $S$ 上光滑, 则:}

\begin{enumerate}[label=(\roman*),leftmargin=2.6em]
\item \underline{层 $R^if_*\Omega^j_{Y/S}$ 局部自由, 且其构造与任意基变换相容.}

\item \underline{联系 Hodge 上同调与 de Rham 上同调的谱序列}
\begin{equation*}
E_1^{p,q}=R^qf_*\Omega^p_{Y/S}\Longrightarrow
R^{p+q}f_*(\Omega^*_{Y/S}),
\end{equation*}
\underline{退化 ($E_1=E_\infty$); 其极限局部自由, 且其构造与任意基变换相容.}
\end{enumerate}

\noindent(iii)$_a$ \underline{有}

\begin{equation*}
R^if_*\Omega^j_{Y/S}=0
\qquad\underline{\text{若 }}i\ne j\ \underline{\text{ 且 }}\ i+j\ne n.
\end{equation*}

%%%%%%%%%% source scan idx 54 | folio 47 %%%%%%%%%%

\noindent(iii)$_b$ \underline{若 $2i<n$, 则 $\eta^i$ 为 $R^if_*\Omega^i_{Y/S}$ 的一组基.}

\noindent(iii)$_c$ \underline{若 $n<2i\leq2n$, 且 $\alpha_i$ 为 $R^if_*\Omega^i_{Y/S}$ 中与 $R^{n-i}f_*\Omega^{n-i}_{Y/S}$ 的基 $\eta^{n-i}$ 对偶的基, 则}

\begin{equation*}
\eta^i=\left(\prod_{h=1}^d a_h\right)\alpha_i.
\end{equation*}

\noindent(iii)$_d$ \underline{若 $n=2m$, 则 $\eta^m\ne0$ (只要 $S\ne\varnothing$), 且商 $R^mf_*\Omega^m_{Y/S}/\underline O_S\eta^m$ 局部自由.}

只须在``普遍''情形下证明 1.5, 这使我们可以假设 $S$ 在 $\operatorname{Spec}(\mathbb Z)$ 上光滑. 此时 (ii) 由 (i) 得到, 因为由 Hodge 理论, (ii) 在 $\mathbb C$ 上成立, 从而在每个零特征域上成立. 此外, 由于 $S$ 为约化概形, 若假定在 $S$ 为域的谱时已证明 (iii), 则 (i) 由 (iii) 得到. 事实上, 由 (iii), 对每个 $j$, 函数

\begin{equation*}
s\longmapsto\dim_{k(s)}H^i(Y_s,\Omega^j_{Y_s})
\end{equation*}

在 $S$ 上均为常数, 至多 $i+j=n$ 时例外\EditorialFootnote{法文原印本与英文底本此处均写作 $i=j$. 由 (iii)$_a$--(iii)$_c$, 对固定的 $j$, 唯一尚待确定的维数项是 $i+j=n$ 的中间上同调项.}; 因其交错和为常数 (EGA III, 7.9.2), 故它们全都为常数. 再由 EGA III, 6.10.5, 7.6.9 和 7.7.5 即得结论.

由 1.4 和 1.3, 有

\begin{equation}\tag{1.5.1}
\left\{
\begin{aligned}
R^ig_*\Omega^j_{\mathbb P(\underline E)/S}
&\xrightarrow{\ \sim\ }R^if_*\Omega^j_{Y/S}
&&\text{当 }i+j<n,\\
R^ig_*\Omega^j_{\mathbb P(\underline E)/S}
&\lhook\joinrel\longrightarrow R^if_*\Omega^j_{Y/S}
&&\text{当 }i+j=n.
\end{aligned}
\right.
\end{equation}

当 $S$ 为域的谱时, 断言 (iii)$_a$, (iii)$_b$ 和 (iii)$_d$ 由 (1.5.1) 与 Serre 对偶性 ([2], V, 11.2, 第~211 页) 得到. 上同调类

%%%%%%%%%% source scan idx 55 | folio 48 %%%%%%%%%%

$\eta^n$ 是次数为 $\prod_{h=1}^d a_h$ (即完全交的次数) 的零圈的类, 所以

\begin{equation*}
\eta^n=\left(\prod_{h=1}^d a_h\right)\alpha_n.
\end{equation*}

根据定义, $\eta^{n-i}\alpha_i=\alpha_n$, 从而得到 (iii)$_c$.

于是 (i) 和 (ii) 成立; 在普遍情形下, (iii) 中各断言由 (i) 和已经证明的 $S$ 为域的谱时的 (iii) 得到.

\underline{定理} 1.6. \underline{保留 1.4 和 1.5 的假设与记号, 假定 $Y$ 光滑, 设 $\ell$ 为在 $S$ 的所有剩余域中可逆的素数, 并以 $\eta_\ell\in H^0(S,R^2f_*\mathbb Z_\ell(1))$ 表示 $\underline O(1)$ 的 $\ell$ 进第一 Chern 类 (SGA 5, VII). 则:}

\begin{enumerate}[label=(\roman*),leftmargin=2.6em]
\item \underline{若 $i$ 为奇数且 $i\ne n$, 则 $R^if_*\mathbb Z_\ell=0$.}

\item \underline{若 $2i<n$, 则 $R^{2i}f_*\mathbb Z_\ell(i)$ 同构于 $\mathbb Z_\ell$, 并以 $\eta_\ell^i$ 为基.}

\item \underline{若 $n<2i\leq2n$, 则 $R^{2i}f_*\mathbb Z_\ell(i)$ 同构于 $\mathbb Z_\ell$, 并以}
\begin{equation*}
\eta_\ell^i\bigg/\prod_{h=1}^d a_h
\end{equation*}
\underline{为基.}

\item \underline{$\mathbb Z_\ell$-层 $R^nf_*\mathbb Z_\ell$ 可构造、扭常且无挠. 若 $n$ 为偶数, 则 $\eta_\ell^{n/2}\ne0$ (只要 $S\ne\varnothing$), 且 $R^nf_*\mathbb Z_\ell(n/2)/\mathbb Z_\ell\eta_\ell^{n/2}$ 无挠.}
\end{enumerate}

由 SGA 4, XVI, 2.2, $\mathbb Z_\ell$-层 $R^if_*\mathbb Z_\ell$ 可构造且扭常, 其构造与任意基变换相容. 因此, 只须在 $S$ 为代数闭域的谱时证明 1.6. 此时反复应用 SGA 4, XIV, 3.3 可知, 对每个 $m$, 典范映射

\begin{equation*}
H_Y^q\bigl(\mathbb P(\underline E),\mathbb Z/\ell^m\bigr)
\longrightarrow H^q(Y,\mathbb Z/\ell^m)
\end{equation*}

%%%%%%%%%% source scan idx 56 | folio 49 %%%%%%%%%%

当 $q>r+d$ 时为双射, 当 $q=r+d$ 时为满射. 由此和 Poincaré 对偶性 (SGA 4, XVIII), 推得 (SGA 5, VII) 限制映射

\begin{equation*}
H^q\bigl(\mathbb P(\underline E),\mathbb Z_\ell\bigr)
\longrightarrow H^q(Y,\mathbb Z_\ell)
\end{equation*}

当 $q<n$ 时为双射, 当 $q=n$ 时为单射且余核无挠. 这证明了 $i<n$ 时的 (i), 以及 (ii) 和 (iv). 与 1.5 中相同, 可验证 $\eta_\ell^n$ 是典范类的 $\prod_{h=1}^d a_h$ 倍; 再由 Poincaré 对偶性 (SGA 4, XVIII) 得到 (i) 和 (iii).

